% Exemplo: construção de tabelas passo a passo: tabular, linhas, colunas de parágrafo, tabularx,
% tabulary, booktabs, mesclagem de células e cores
% Compilar: latexmk -pdf tabular-basico.tex        Testado com: TeX Live 2024 ou mais recente
% Cada trecho fica em uma página, para que o resultado de cada um possa ser mostrado separadamente.
% Para tabelas no padrão ABNT/IBGE (título, fonte e nota), veja tabela-simples.tex.
\documentclass[a4paper]{article}
\usepackage[T1]{fontenc}
\usepackage{lmodern}
\usepackage[brazilian]{babel}
%%BEGIN:tabelaspreambulo
\usepackage{array}             % colunas m{} e b{}
\usepackage{tabularx}          % colunas X (largura total fixa)
\usepackage{tabulary}          % colunas L, C, R, J (largura máxima)
\usepackage{booktabs}          % \toprule, \midrule, \bottomrule
\usepackage{multirow}          % \multirow
\usepackage[table]{xcolor}     % cores em tabelas
%%END:tabelaspreambulo

\begin{document}

%%BEGIN:tabularsimples
\begin{tabular}{ccl}
  \textbf{Qtd} & \textbf{Medida}  & \textbf{Ingrediente} \\
  500          & g                & polvilho azedo \\
  1            & copo americano   & água \\
  1            & copo americano   & leite \\
  1/2          & xícara           & óleo \\
  2            & un               & ovos \\
  100          & g                & queijo parmesão ralado \\
  ---          & ---              & sal a gosto \\
\end{tabular}
%%END:tabularsimples

\clearpage

%%BEGIN:tabularlinhas
\begin{tabular}{|c|c|l|}
  \hline
  \textbf{Qtd} & \textbf{Medida} & \textbf{Ingrediente} \\
  \hline
  500 & g              & polvilho azedo \\
  \cline{2-3}
  1   & copo americano & água \\
  \cline{1-2}
  1   & copo americano & leite \\ \hline
  1/2 & xícara         & óleo \\ \hline
\end{tabular}
%%END:tabularlinhas

\clearpage

%%BEGIN:colunalonga
\begin{tabular}{|l|c|l|}
  \hline
  1 & 2 & 3 \\ \hline
  4 & 5 & Célula com muito texto, que alarga a coluna até sair da margem da página \\ \hline
  7 & 8 & 9 \\ \hline
\end{tabular}
%%END:colunalonga

\clearpage

%%BEGIN:colunap
\begin{tabular}{|l|c|p{4cm}|}
  \hline
  1 & 2 & 3 \\ \hline
  4 & 5 & Célula com muito texto, que agora quebra em várias linhas \\ \hline
  7 & 8 & 9 \\ \hline
\end{tabular}
%%END:colunap

\clearpage

%%BEGIN:tabularx
\begin{tabularx}{.6\linewidth}{|l|c|X|}
  \hline
  1 & 2 & 3 \\ \hline
  4 & 5 & Célula com muito texto: a coluna X ocupa a largura que sobra \\ \hline
  7 & 8 & 9 \\ \hline
\end{tabularx}
%%END:tabularx

\clearpage

%%BEGIN:colunasxiguais
\begin{tabularx}{\linewidth}{X|X|X}
  \hline
  Esticar esta coluna & 2 & 3 \\ \hline
  4 & 5 & Esticar esta coluna \\ \hline
  7 & Esta também & 9 \\ \hline
\end{tabularx}
%%END:colunasxiguais

\clearpage

%%BEGIN:tabulary
\begin{tabulary}{\linewidth}{|L|C|L|}
  \hline
  Esticar esta coluna & 2 & 3 \\ \hline
  4 & 5 & Esticar esta coluna \\ \hline
  7 & Esta também & 9 \\ \hline
\end{tabulary}
%%END:tabulary

\clearpage

%%BEGIN:booktabs
\begin{tabular}{ccl}
  \toprule
  \textbf{Qtd} & \textbf{Medida} & \textbf{Ingrediente} \\
  \midrule
  500 & g              & polvilho azedo \\
  1   & copo americano & água \\
  1   & copo americano & leite \\
  1/2 & xícara         & óleo \\
  2   & un             & ovos \\
  100 & g              & queijo parmesão ralado \\
  --- & ---            & sal a gosto \\
  \bottomrule
\end{tabular}
%%END:booktabs

\clearpage

%%BEGIN:mesclagemcolunas
\begin{tabular}{llr}
  \toprule
  \multicolumn{3}{c}{\textbf{Matriz curricular (trecho)}} \\
  \midrule
  1º período & Cálculo I                 & 60 \\
  1º período & Introdução à Programação  & 60 \\
  1º período & Metodologia Científica    & 30 \\
  2º período & Cálculo II                & 60 \\
  2º período & Estruturas de Dados       & 60 \\
  \bottomrule
\end{tabular}
%%END:mesclagemcolunas

\clearpage

%%BEGIN:mesclagemlinhas
\begin{tabular}{|l|l|r|}
  \hline
  \multicolumn{3}{|c|}{\textbf{Matriz curricular (trecho)}} \\ \hline
  \multirow{3}{*}{1º período} & Cálculo I                & 60 \\
                              & Introdução à Programação & 60 \\
                              & Metodologia Científica   & 30 \\ \hline
  \multirow{2}{*}{2º período} & Cálculo II               & 60 \\
                              & Estruturas de Dados      & 60 \\ \hline
\end{tabular}
%%END:mesclagemlinhas

\clearpage

%%BEGIN:coresalternadas
\rowcolors{2}{gray!15}{white}   % a partir da linha 2: cinza nas pares, branco nas ímpares
\begin{tabular}{lr}
  \textbf{Disciplina} & \textbf{Média} \\
  Cálculo I           & 7,2 \\
  Programação         & 8,1 \\
  \hiderowcolors                  % desliga as cores a partir daqui
  Metodologia         & 8,8 \\
  \showrowcolors                  % religa as cores
  Estruturas de Dados & 6,9 \\
  Bancos de Dados     & 7,5 \\
\end{tabular}
%%END:coresalternadas

\clearpage

%%BEGIN:corescelulas
\begin{tabular}{lr}
  \rowcolor{gray!30}
  \textbf{Disciplina} & \textbf{Média} \\
  Cálculo I           & 7,2 \\
  Estruturas de Dados & \cellcolor{red!20}5,9 \\
  Programação         & \cellcolor{green!20}8,1 \\
\end{tabular}
%%END:corescelulas

\end{document}
